El ${}^{14}_{6}\mathrm{C}$ es produeix a l'atmosfera per l'acció dels raigs còsmics. Aquest isòtop és inestable i va decaient a ${}^{14}_{7}\mathrm{N}$ mitjançant un procés de desintegració $\beta$, amb un període de semidesintegració de 5\,730 anys. La proporció de ${}^{14}\mathrm{C}$ respecte al ${}^{12}\mathrm{C}$ que hi ha a l'atmosfera és constant al llarg del temps. Els éssers vius assimilen el CO$_2$ de l'atmosfera sense distingir si es tracta de ${}^{12}\mathrm{C}$ o de ${}^{14}\mathrm{C}$, i ho fan en la proporció en què aquests isòtops es troben de manera natural a l'atmosfera. Quan moren, els éssers deixen d'assimilar CO$_2$ i, a partir d'aquest moment, la quantitat de ${}^{14}\mathrm{C}$ va decaient.

\begin{apartat}{1,25}
Escriviu la reacció que correspon al decaïment del ${}^{14}\mathrm{C}$ a ${}^{14}\mathrm{N}$. Incloeu-hi, si escau, els antineutrins.
\end{apartat}

\begin{apartat}{1,25}
Si una mostra d'una fusta utilitzada en un sarcòfag presenta una proporció de ${}^{14}_{6}\mathrm{C}$ de només el 58\,\% respecte a la proporció que hi ha a l'atmosfera, trobeu quina és l'antiguitat del sarcòfag.
\end{apartat}
